
\section{Algebraic Number Theory}
\subsection{Why New Numbers Appear}

The natural numbers are sufficient to solve equations such as

\[
x+3=5.
\]

However, they are not sufficient to solve

\[
x+5=3.
\]

This leads to the introduction of negative numbers.

Likewise, integers are not sufficient to solve

\[
2x=1,
\]

which motivates the rational numbers.

The rational numbers are not sufficient to solve

\[
x^2=2,
\]

which motivates the real numbers.

The real numbers are not sufficient to solve

\[
x^2+1=0,
\]

which motivates the complex numbers.

A recurring pattern appears throughout mathematics:

\begin{quote}
When an equation has no solution in the current number system, we enlarge the system.
\end{quote}

The study of numbers is therefore inseparable from the study of equations.

\section{A New Number}

Consider the equation

\[
x^2+5=0.
\]

Introduce a symbol (u) satisfying

\[
u^2=-5.
\]

We denote this new number system by

\[
\mathbb Z[u]
=
\{a+bu : a,b\in\mathbb Z\}.
\]

Since \(u^2=-5\), we may also write

\[
\mathbb Z[\sqrt{-5}].
\]


We now allow expressions of the form

\[
a+bu,
\]

where $(a,b\in\mathbb Z)$.

Examples are

\[
3+u,\qquad 7-2u,\qquad 5.
\]

This creates a new arithmetic world containing the integers together with a solution of

\[
x^2+5=0.
\]

\section{An Unexpected Observation}

In the integers, the number

\[
6
\]

has the factorization

\[
6=2\cdot3.
\]


Using $(u^2=-5)$, we compute

\[
(1-u)(1+u)=1-u^2=1-(-5)=6.
\]

Therefore

\[
6=(1-u)(1+u).
\]

At first sight this appears harmless.

Is this really a different factorization of \(6\)?

The next problems investigate this question.

This raises a fundamental question:

\begin{quote}
Can arithmetic behave differently in different number systems?
\end{quote}

The remainder of this book is an attempt to understand this phenomenon.

\subsection{Problem 1}
Compute: 
\[(1+u)(1-u)\]
Use the fact:
\[u^2 = -5\]

What ordinary integer do you obtain?

\subsection{Problem 2}
Compute:
\[(2+u)(2-u)\]

\subsection{Problem 3}
Look the number:
\[6 = 3.2\]

Now compare with your answers from problems 1 and 2.

\subsection{Problem 4}
Can 2 divide 1+u such that:
 \[1+u = 2(a+bu)\]?

 \subsection{Problem 5}
Can 3 divide 1+u such that:
 \[1+u = 3(a+bu)\]?

Surprisingly, none of the factors appears to divide the others.

This suggests that

\[
6 = 2\cdot3=(1-u)(1+u)
\]

may genuinely be two different factorizations of the same number.

If this is true, then unique factorization has failed.

\section{Norm}

The factorization

\[
(a+bu)(a-bu)
=
a^2+5b^2
\]

produces an ordinary integer.

This suggests associating to each number \(a+bu\) the integer

\[
N(a+bu)=a^2+5b^2,
\]
a ordinary integer, that we will call $\textit{Norm}$ and use the symbol $\textit{N}$ to compute it.
The norm allows us to measure elements of \(\mathbb Z[u]\) using ordinary integers.

Since

\[
N(a+bu)
=
(a+bu)(a-bu),
\]

many questions about factorization in \(\mathbb Z[u]\) can be translated into questions about ordinary integers.
\subsection{Problem 6}

Prove that

\[
N(\alpha\beta)
=
N(\alpha)N(\beta).
\]

This property is called multiplicativity of the norm.

\subsection{Problem 7}
Compute the norms:
\begin{align*}
N(1+u) \\
N(2) \\
N(3)
\end{align*}

Now suppose

\[
1+u=\alpha\beta.
\]

Using the multiplicativity of the norm,

\[
N(1+u)
=
N(\alpha)N(\beta).
\]

Since

\[
N(1+u)
=
1^2+5(1)^2
=
6,
\]

we obtain

\[
6=N(\alpha)N(\beta).
\]

The positive factorizations of \(6\) are

\[
1\cdot6
\]

and

\[
2\cdot3.
\]

Can any element of \(\mathbb Z[u]\) have norm \(2\)?

Can any element of \(\mathbb Z[u]\) have norm \(3\)?

If not, the only possibility is

\[
N(\alpha)=1
\]

or

\[
N(\beta)=1.
\]

A norm equal to \(1\) corresponds to a unit.

Therefore one factor must be a unit, and \(1+u\) is irreducible.

\begin{quote}
    We said irreducible not prime.
\end{quote}


\subsection{Problem 8}
In ordinary integers, if $\textit{p}$ is prime and
    \[p|ab\]

then
    \[p|a\]

or

    \[p|b\]

Now test this here with:
    \[3|(1+u)(1-u)\]

Why does 3 divide the product but not the either factor?

\begin{quote}
    Answering this makes you start understanding why we do not use the prime word rather 
    irreducible.
\end{quote}

    \begin{align*}
        1. \text{What should irreducible mean?} \\
        2. \text{What should prime mean?  }      
    \end{align*}

\section{irreducible and Primes}

\subsection{irreducibles}

From our observations we can define that a number \textit{r} is irreducible if whenever
    \[r=ab\]
one of the factors is a unit, 1 or -1.
In the world of a+bu the only units are 1 and -1.

Thus 1+u is irreducible because every factorization has unit factor.

\subsection{Primes}

Now for prime using our new world $\mathbb{Z}[\sqrt{-5}]$
we already discover that 3 fails because it is irreducible but not
prime.

Remember that a number $\textit{p}$ is prime if whenever:
    \[p|ab\]

then:
    \[p|a\]

or \[p|b\]

Notice that this definition does not mention factorization of p itself. It is a divisibilty
property.

Now historically in ordinary integers:
    \begin{quote}
        irreducible = prime
    \end{quote}

But in our new world $\mathbb{Z}[\sqrt{-5}]$ they are different.

\subsection{Problem 9}
Can you prove that 1+u, 1-u and 2 are irreducible?

Hint: Use the norm N.


\section{Where We Are}

We began by enlarging the integers in order to solve an equation.

The new numbers behaved naturally under addition and multiplication.

Yet something unexpected occurred: a familiar theorem about unique factorization appeared to fail.

The number

\[
6 = 2\cdot3 = (1-u)(1+u)
\]

seems to admit two genuinely different factorizations.

The remainder of this book investigates how arithmetic changes when the number system changes.

\vspace{0.5cm}

What should we do next?

\begin{itemize}
\item Redefine primes?
\item Redefine factorization?
\item Search for a new object that factors uniquely even when numbers do not?
\end{itemize}

The third possibility turns out to be the most fruitful.

After all, the numbers themselves seem perfectly reasonable.

Addition works.

Multiplication works.

The norm works.

What appears to have failed is factorization.

This suggests a new question.

\begin{quote}
Is there a hidden object behind numbers that still factors uniquely?
\end{quote}

Let us return to the mysterious identity

\[
6=2\cdot3=(1+u)(1-u).
\]

Imagine there were objects

\[
A,\quad B,\quad C,\quad D
\]

such that

\[
2=AB,
\]

\[
3=CD,
\]

and

\[
1+u=AC,
\]

\[
1-u=BD.
\]

Then we would obtain

\[
6 = 2\cdot3 = (AB)(CD) =
ABCD
\]

and also

\[
6 = (1+u)(1-u) = (AC)(BD) = ABCD.
\]

Suddenly the factorization becomes unique again.

We do not know whether such objects actually exist.

We only know what properties they would need to have.

This gives us a clue.

\begin{quote}
Perhaps numbers are only shadows of a finer factorization system.
\end{quote}

The next step is to investigate divisibility more carefully.

\section{Divisibility Sets}

Consider all numbers divisible by (2).

In ordinary integers this gives

\[
{\ldots,-4,-2,0,2,4,\ldots}.
\]

What does the analogous set look like in $\mathbb Z[\sqrt{-5}]$?

\subsection{Problem 10}

Assume a+bu is divisible by 2.

Then

\[
a+bu=2(c+du).
\]

What conditions must (a) and (b) satisfy?

\subsection{Problem 11}

Describe explicitly the set of all numbers divisible by 2.

Call this set \{2\},

\subsection{Problem 12}

Suppose

\[
x,y\in\{2\}.
\]

Show that

\[
x+y\in\{2\}.
\]

In other words, the sum of two elements divisible by 2 is again divisible by 2.

\subsection{Problem 13}

Suppose

\[
x\in\{2\}
\]

and

\[
r\in\mathbb Z[\sqrt{-5}].
\]

Show that

\[
rx\in\{2\}.
\]

Thus multiplying by any element of the $\mathbb Z[\sqrt{-5}]$ keeps us inside the set.

We have discovered a remarkable kind of subset.

It is closed under addition.

It is also stable under multiplication by arbitrary elements of the $\mathbb Z[\sqrt{-5}]$.

Such subsets appear naturally whenever we study divisibility.

Because they occur so frequently, they deserve a name.

\section{Ideals}

An $\textit{ideal}$ is a subset I of ($\mathbb Z[\sqrt{-5}]$) satisfying:

\begin{enumerate}
\item If x,y$\in$ I, then x+y$\in$ I.
\item If x$\in$ I and r$\in\mathbb Z[\sqrt{-5}$], then rx$\in $I.
\end{enumerate}

The set:

\[
\{2\} = \{2(a+bu):a,b\in\mathbb Z\}
\]

is our first example.

Likewise we may consider

\[
\{3\},
\]

the set of all numbers divisible by 3,

and

\[
\{(1+u)\},
\]

the set of all numbers divisible by (1+u).

At first sight these sets seem to contain less information than the numbers themselves.

Surprisingly, the opposite is true.

As we shall discover, ideals remember factorization more faithfully than numbers do.

\section{Prime Ideals}

Now to save factorization again we start thinking what would be a prime ideal?
At the beginning of mathematics, numbers seemed fundamental but after enough algebraic
number theory, one starts thinking:

    \begin{quote}
        The polynomial is primary and the numbers are shadows of the polynomial.
    \end{quote}

For example:
    \[x^2+1\]

creates the Gaussian integers.

    \[x^2+5\]

creates $\mathbb{Z}[\sqrt{-5}]$.

    \[x^2-2\]

creates $\mathbb{Z}[\sqrt{2}]$.

Different polynomials create a different arithmetic universes.

\subsection{Problem 14}
Suppose $\alpha$ satisfies 
    \[\alpha^2-2 = 0\]

abd $\beta$ satifies
    \[\beta^4-4\beta^2+2 = 0\]

How could you tell whether $\beta$ lives in the same arithmetic world as $\alpha$?

I hope you did a break and try to solve it because now you are in the edge of a great discover.
Treat y = $x^2$:
    \[y^2-4y+2\]

Solving Baskara we find:
    \[y = \frac{4\pm\sqrt{16-8}}{2} = 2\pm\sqrt{2}\]

Therefore the roots satisfy:
    \[x^2+\sqrt{2}\]
or
    \[x^2-\sqrt{2}\]

Now, notice something remarkable. The polynomial contains only integers:
    \[x^4-4x^2+2\]

Yet, it's roots invole $\sqrt{2}$, so:
    \begin{quote}
        The world generated by it's roots seems connected to the world generated by $\sqrt{2}$.
    \end{quote}

This makes us think:
    \begin{quote}
        Does a root of one polynomial already live inside the arithmetic universe generated by another root?
    \end{quote}

or

    \begin{quote}
        When do two different polynomials generate the same arithmetic world?
    \end{quote}